\subsection{Ampliação de escala (\textit{scale-up})}

Tendo sido a amostra do quinto dia o melhor produto obtido, sua formulação foi tomada como base a ampliação de escala (bancada → industrial). Como o procedimento empregado é uma saponificação direta a quente, seu análogo industrial mais direto é o processo de caldeira (\textit{kettle process}), descrito na Seção~\ref{subsubsec:kettle}.

\subsubsection{Balanço de massa}

A ampliação é definida por um fator de escala $S$, dado pela razão entre a massa de óleo da batelada industrial pretendida e a da bancada (Equação~\ref{eq:scaleup_fator}). Tomando como base uma batelada industrial de uma tonelada de óleo, tem-se $S = 1000~\text{kg} / 128{,}79~\text{g} \approx 7{,}8 \times 10^{3}$.

\begin{equation}
    S = \frac{M_{\text{óleo}}^{\text{ind}}}{M_{\text{óleo}}^{\text{banc}}}
    \label{eq:scaleup_fator}
\end{equation}

Como a formulação é mantida, a massa de cada componente escala linearmente com $S$. A Tabela~\ref{tab:scaleup} apresenta a formulação do quinto dia e a correspondente para uma batelada industrial de uma tonelada de óleo.

\begin{table}[H]
    \centering
    \footnotesize
    \caption{Formulação do quinto dia escalada para uma batelada industrial de uma tonelada de óleo.}
    \label{tab:scaleup}
    \begin{tabular}{lcc}
        \toprule
        \textbf{Componente} & \textbf{Dia 5} & \textbf{Industrial (por tonelada de óleo)} \\
        \midrule
        Óleo de babaçu     & 45,72~g   & 355~kg \\
        Óleo de soja       & 83,07~g   & 645~kg \\
        Ácido esteárico    & 4,05~g    & 31,5~kg \\
        NaOH               & 19,18~g   & 149~kg \\
        Água               & 32,2~ml   & 250~L \\
        Etanol (95\%)      & 7,30~ml   & 56,7~L \\
        \bottomrule
    \end{tabular}
\end{table}

\subsubsection{Coproduto e adaptações de processo}

Diferentemente da bancada, que retém o glicerol no produto, o processo de caldeira o separa por \textit{salting-out}. A massa de glicerol gerada é estequiométrica e da ordem de 10\% da massa de óleo, o que corresponde a cerca de 100~kg por tonelada de óleo processada. Esse seria um coproduto de valor agregado, que pode ser recuperado e, se desejado, parcialmente adicionado de volta ao sabão na etapa de acabamento, conforme discutido anteriormente.

A transferência para a escala industrial exige ainda adaptar as condições físicas do processo. O aquecimento por chapa e a agitação por agitador de bancada dão lugar à injeção direta de vapor e à circulação nas caldeiras. Introduz-se a etapa de \textit{salting-out}, responsável tanto pela recuperação do glicerol quanto pela purificação do sabão (remoção do excesso de álcali, do sal e da matéria não saponificada), a moldagem e a cura manuais são substituídas por linhas de secagem, extrusão e estampagem \cite{shreve1984chemical}. O etanol, empregado na bancada como cossolvente, em geral não é usado em escala industrial, dado que o controle de homogeneização passa a ser feito pela própria fervura e seria um material caro nessa escala. Por fim, o controle do fim da reação deixa de ser feito por inspeção visual e pH e passa a ser analítico (álcali livre, sal e umidade), e o ácido esteárico é incorporado na etapa de acabamento.
